% answers_short.tex —— 2010 年·普通高等学校招生考试·安徽卷·文科 · 简短答案速查
% 由 tools/make_short_answers.py 自 answers.tex 的 \gksolhead/\gkanswer 提取，
% 并与 2010/2010_数学真题_安徽文_答案要点.md 逐题比对后生成，勿手改（改源 MD 后重跑）。

\begin{gknotice}
\textbf{使用说明}：本节只给最终答案便于快速核对，完整推导见解析版。答案为编写者独立推导，\textbf{非官方标准答案}。
\end{gknotice}

\gkpart{选择题}
\begin{center}\begin{tabular}{ccccccccccc}
\toprule
题号 & 1 & 2 & 3 & 4 & 5 & 6 & 7 & 8 & 9 & 10 \\
\midrule
答案 & C & B & D & A & A & D & A & C & B & C \\
\bottomrule
\end{tabular}\end{center}

\gkpart{填空题}
\begin{itemize}[leftmargin=2.6em]
  \item[\textbf{11.}] 对任意 $x\in\mathbf R$，都有 $x^2+2x+5\neq0$。
  \item[\textbf{12.}] $(1,0)$。
  \item[\textbf{13.}] $10$。
  \item[\textbf{14.}] $11\%$（即 $\dfrac{12}{109}$）。
  \item[\textbf{15.}] ①③⑤。
\end{itemize}

\gkpart{解答题}
\begin{itemize}[leftmargin=2.6em]
  \item[\textbf{16.}] （1）$144$；（2）$5$。
  \item[\textbf{17.}] （1）$\dfrac{x^2}{16}+\dfrac{y^2}{12}=1$；（2）$2x-y-1=0$。
  \item[\textbf{18.}] （1）频率分布表如下；（2）直方图如下；（3）空气质量总体为良。
  \item[\textbf{19.}] （1）证明；（2）证明；（3）$\dfrac13$。
  \item[\textbf{20.}] 增区间为 $(0,\pi)$ 与 $\left(\dfrac{3\pi}2,2\pi\right)$，减区间为 $\left(\pi,\dfrac{3\pi}2\right)$；极大值 $\pi+2$，极小值 $\dfrac{3\pi}2$。
  \item[\textbf{21.}] （1）证明；（2）$\dfrac94-\dfrac{2n+3}{4\cdot3^{n-1}}$。
\end{itemize}
